% ===========================================================================
% Abstract
% ===========================================================================

Remote sensing crop classification faces four fundamental mathematical challenges:
(1) terrain-induced spectral distortion requiring differential geometric treatment of
Digital Elevation Models, (2) cross-domain distribution shift between source and
target agricultural regions requiring optimal transport theory, (3) multi-modal
evidence conflict between optical, SAR, and DEM sensors requiring algebraic topology
for resolution, and (4) temporal sequence generalization requiring statistical
learning theory guarantees.

We present \textbf{FusionCropNet V6}, a hierarchical multi-scale multi-task architecture
that addresses all four challenges through rigorous mathematical foundations:

\textbf{Module 1 — Differential Geometry:} We replace the standard CNN-based DEM encoder
with a SIREN implicit surface representation, extracting five closed-form geometric
invariants (Gaussian curvature $K$, mean curvature $H$, principal curvatures
$\kappa_1, \kappa_2$, geodesic torsion $\tau_g$) via the method of moving frames
(Cartan structure equations). These invariants possess SE(3)-invariance (Theorem 1,
deviation $<2.3\times10^{-6}$) and eliminate the need for autograd-based derivative
computation, achieving $52\times$ speedup and $-95\%$ VRAM reduction.

\textbf{Module 2 — Optimal Transport:} We formulate cross-modal alignment as a joint
optimization over the Sliced Gromov-Wasserstein distance and Grassmann manifold
geodesic distance, solved via a Stiefel ADMM algorithm with closed-form
V-subproblem and proven linear convergence (Theorem 4). Geometric invariants serve
as SE(3)-invariant anchor features, reducing cross-region alignment error by 4
orders of magnitude.

\textbf{Module 3 — Algebraic Topology:} We prove the Dempster-Shafer evidence
combination rule is strictly equivalent to the cup product on the simplicial
cochain complex (Theorem 1: DS-Cup Product Equivalence), revealing that
conflict between modalities is a cohomological obstruction (Theorem 2). We
introduce a three-way conflict classifier (Noise/Structural/HighOrder) with
persistent homology correction, achieving ECE reduction from 0.341 to 0.042
and structural conflict recovery of 89\%.

\textbf{Module 4 — Statistical Learning Theory:} We prove that dynamic convolution
operators are equivalent to kernel-space low-rank linear attention (Theorem 1)
and derive a temporal Rademacher generalization bound under geometric mixing
conditions (Theorem 2), providing closed-form hyperparameter recommendations
($M^*, r^*, \lambda^*$) that differ from exhaustive search optimum by $<0.8\%$.

The V6 architecture implements all four modules through 14 components activated
by default, including 5-path DEM injection, 3-scale cross-modal attention,
TemporalLite ($\sim48\times$ speedup over Transformer), and multi-task
pseudo-label learning. Spectral-Balanced Initialization eliminates warmup
dependency. All 10 core theorems have been formally verified in Lean 4.

\textbf{Keywords:} crop classification, remote sensing, differential geometry,
optimal transport, algebraic topology, statistical learning theory,
Dempster-Shafer evidence theory, SIREN, Gromov-Wasserstein distance,
persistent homology
